% 1982 年全国卷（文科）第 5 题：以墙为一边、用篱笆围长方形场地，并以平行于宽边的篱笆隔成两块；上沿为带砖缝与短刻线的墙，墙两端略长于场地，左块窄、右块宽（按源卷布局）
\begin{tikzpicture}[scale=1, line width=0.5pt,
  every node/.style={inner sep=1.5pt}]
  % 墙：上沿 y=4.00、下沿 y=3.60 两条水平线
  \draw (0,4.00) -- (8,4.00);
  \draw (0,3.60) -- (8,3.60);
  % 砖缝（两线之间）
  \foreach \x in {0.65,1.45,2.25,3.05,3.85,4.65,5.45,6.25,7.05,7.85} {
    \draw (\x,3.60) -- (\x,4.00);
  }
  % 上沿短刻线
  \foreach \x in {0.25,1.05,1.85,2.65,3.45,4.25,5.05,5.85,6.65,7.45} {
    \draw (\x,4.00) -- (\x,4.27);
  }
  % 下沿短刻线
  \foreach \x in {0.25,1.05,1.85,2.65,3.45,4.25,5.05,5.85,6.65,7.45} {
    \draw (\x,3.60) -- (\x,3.33);
  }
  % 场地：左篱、中隔篱（平行于两侧的篱笆）、右篱均上端接墙的下沿，另加下横篱
  \draw (1.00,3.60) -- (1.00,0.38);
  \draw (2.95,3.60) -- (2.95,0.38);
  \draw (6.90,3.60) -- (6.90,0.38);
  \draw (1.00,0.38) -- (6.90,0.38);
\end{tikzpicture}
